% ===================================================================
%  اللوحة رقم ١٢ — المحطة. — السكك الحديدية. — السفن.
%  Traduction 2026 d'apres texto/io, controlee sur texto/fr, texto/en
%  et texto/es. Consigne : voir texto/ar/10-lawha-01.tex.
% ===================================================================

%%K t12-apar-1 apar
\VUsaut{4.0mm}
\VUtitre{15.0pt}{60}{اللوحة رقم ١٢}
\VUsaut{2.4mm}
\VUpk{11.4pt}{40}{المحطة. --- السكك الحديدية. --- السفن.}
\VUsaut{3.4mm}
\textsc{ملاحظة.} --- \textit{المواقيت المستعملة في كل بلد مختلفة، ولذلك
اتخذنا ساعات} \VUgras{اللوحة الفرنسية} \textit{مثالًا.}

%%K t12-01-1 p
1. --- قمت لتوّي برحلة في ألمانيا. وقبل السفر كنت قد اشتريت \VUgras{جدول
مواقيت} \textsuperscript{(15)} لأعرف ساعات قيام القطارات. ولما تبيّن لي أن
أنسب قطار يغادر باريس في التاسعة مساءً ويبلغ ستراسبورغ في الواحدة وثلاث
عشرة دقيقة، تأهبت للسفر، وفي الغد أمرت بأن أُحمَل إلى المحطة.

%%K t12-02-1 p
2. --- ولما دخلت بهو القيام الكبير خلف \VUgras{كاهن} \textsuperscript{(2)} ريفي
عجوز يحمل \VUgras{حقيبة سفره} \textsuperscript{(3)}، كان عدد غير قليل من
\VUgras{المسافرين} \textsuperscript{(1)} قد وصلوا.

%%K t12-02-2 p
ولأني أردت أن أرسل \VUgras{برقية} \textsuperscript{(5)}، طلبت من \VUgras{مفتّش}
\textsuperscript{(6)} أن يدلّني على \VUgras{مكتب البرق} \textsuperscript{(4)}.

%%K t12-03-1 p
3. --- وقبل أن أنتقل إلى \VUgras{قاعة الانتظار} \textsuperscript{(7)}، ذهبت
لآخذ \VUgras{تذكرة} \textsuperscript{(9)} ذهاب وإياب.

%%K t12-04-1 p
4. --- ولم أنتظر طويلًا عند \VUgras{شباك التذاكر} \textsuperscript{(8)}، فقد كان
الناس قليلين. وتفضّل \VUgras{جندي} \textsuperscript{(10)} ذاهب في إجازة فأعطاني
دوره، فوجدت من الوقت ما يكفي لأسأل \VUgras{ناظر المحطة} \textsuperscript{(13)}
بضعة أسئلة، وكان يتحدث مع \VUgras{مأمور الشرطة} \textsuperscript{(12)}
و\VUgras{الدركيّ} \textsuperscript{(11)} المناوب.

%%K t12-tit-1 sub
\VUpk{10.2pt}{40}{تسجيل الأمتعة.}
\VUsaut{3.0mm}

%%K t12-05-1 p
5. --- وبعد أن اشتريت جريدة من \VUgras{كشك الكتب} \textsuperscript{(14)}، أمرت
بوزن \VUgras{أمتعتي} \textsuperscript{(17)} --- وكان \VUgras{حمّال} \textsuperscript{(16)}
قد جاء بها لتوّه على \VUgras{عربة أمتعة} \textsuperscript{(19)}. وقال لي
\VUgras{الموظف} \textsuperscript{(22)} المكلّف بـ\VUgras{الميزان} \textsuperscript{(21)}
إن صندوقي يزن خمسة وثلاثين كيلوغرامًا، أي بـ\VUgras{زيادة} خمسة
كيلوغرامات، دفعت عنها فرقًا في \VUgras{شباك الأمتعة} \textsuperscript{(23)}.

%%K t12-06-1 p
6. --- ولأني جئت مبكرًا، وجدت وقتًا لأرقب حركة المسافرين في المحطة. فمنهم
من يترك حقائب صغيرة أو يستردّها في \VUgras{مكتب الأمانات} \textsuperscript{(25)}؛
ومنهم من يراجع \VUgras{جداول المواقيت} \textsuperscript{(26)}. والقادمون يسرعون
إلى \VUgras{المخرج} \textsuperscript{(32)} و\VUgras{حقائبهم} \textsuperscript{(24)} في
أيديهم. وبعضهم ينتظر عند \VUgras{تسليم الأمتعة} \textsuperscript{(27)} ويسلّم
\VUgras{إيصاله} \textsuperscript{(29)} ليأخذ صناديقه أو \VUgras{أقفاصه}
\textsuperscript{(28)} أو \VUgras{سلاله} \textsuperscript{(30)}.

%%K t12-07-1 p
7. --- و\VUgras{موظفو المكوس} \textsuperscript{(33)} يفتشون الطرود بعناية ليروا
أفيها ما يُصرَّح به.

%%K t12-08-1 p
8. --- واتجه بعض المسافرين إلى \VUgras{قاعة المرطبات} \textsuperscript{(34)} حيث
كان \VUgras{نادل} \textsuperscript{(35)} مشغولًا بخدمتهم. وكان عليهم أن يعجّلوا،
فـ\VUgras{القطار} \textsuperscript{(36)}، وهو سريع، لا يقف في المحطة طويلًا.

%%K t12-09-1 p
9. --- ثم خرجت إلى رصيف القيام وبحثت عن \VUgras{عربة} \textsuperscript{(37)}
مباشرة لئلا أضطر إلى التبديل في الطريق. فركبت عربة ذات ممر فيها
\VUgras{مقصورة} \textsuperscript{(38)} للمدخنين وأخرى محجوزة «للسيدات المسافرات
وحدهن».

%%K t12-tit-2 sub
\VUpk{10.2pt}{40}{القيام ليلًا.}
\VUsaut{3.0mm}

%%K t12-10-1 p
10. --- وبعد أن صعدت \VUgras{الدرجة} \textsuperscript{(40)}، وأغلقت \VUgras{الباب}
\textsuperscript{(39)}، ورفعت \VUgras{الستارة} \textsuperscript{(41)}، ووضعت حقائبي
الصغيرة على \VUgras{الرفّ} \textsuperscript{(43)}، جلست على \VUgras{المقعد}
\textsuperscript{(42)}. ولما رأيت \VUgras{عامل المصابيح} \textsuperscript{(44)} يشعل
المصابيح تذكّرت أن أمامي ليلة أقضيها في العربة، فناديت الخادم المكلّف
بتأجير \VUgras{الوسائد} \textsuperscript{(45)} لأطلب منه واحدة.

%%K t12-11-1 p
11. --- وجاء الكمساري ليدقق التذاكر. فسألته لِمَ لم نتحرك وقد حان الوقت.
فأجاب أن \VUgras{رئيس القطار} \textsuperscript{(46)} مشغول بوضع الدراجات في
\VUgras{عربة الأمتعة} \textsuperscript{(47)}. ثم إن \VUgras{مدفآت القدمين}
\textsuperscript{(48)} لم تُبدَّل كلها بعد.

%%K t12-12-1 p
12. --- لكننا كنا على وشك القيام، فـ\VUgras{الوقّاد} \textsuperscript{(50)}، فوق
\VUgras{مقطورة الفحم} \textsuperscript{(49)}، كان يأخذ الفحم ليذكي نار
\VUgras{القاطرة} \textsuperscript{(54)}، و\VUgras{السائق} \textsuperscript{(51)} لا ينتظر
إلا إشارة \VUgras{مساعد ناظر المحطة} \textsuperscript{(52)} الواقف على
\VUgras{الرصيف} \textsuperscript{(53)}.

%%K t12-13-1 p
13. --- و\VUgras{مِرجل} \textsuperscript{(56)} القاطرة يهدر هديرًا خافتًا؛
و\VUgras{المِدخنة} \textsuperscript{(55)} تقذف سيولًا من الدخان مخلوطًا
بـ\VUgras{البخار} \textsuperscript{(60)}. ودوّى \VUgras{صفير} \textsuperscript{(59)} حادّ
وأخذت \VUgras{المكابس} \textsuperscript{(58)} تتحرك مديرةً \VUgras{عجلات}
\textsuperscript{(57)} الآلة الثقيلة.

%%K t12-tit-3 sub
\VUsaut{3.0mm}
\VUpk{10.2pt}{40}{انطلقنا!}
\VUsaut{3.0mm}

%%K t12-14-1 p
14. --- وأخذ القطار، وهو يجري على \VUgras{القضبان} \textsuperscript{(64)}، يزيد
سرعته؛ وانتهت \VUgras{الأرصفة} \textsuperscript{(65)}؛ ومررنا
بـ\VUgras{المراحيض} \textsuperscript{(66)}، وبصفّ من العربات واقف على
\VUgras{خط} \textsuperscript{(63)} آخر: \VUgras{عربات نوم} \textsuperscript{(67)}،
و\VUgras{عربة مطعم} \textsuperscript{(68)}، و\VUgras{عربة بريد} \textsuperscript{(69)}.

%%K t12-noto-1 noto
\VUnotes{143.2mm}{%
\textit{(*) ملاحظة.} --- \textit{الرحلة موصوفة بطبيعة الحال كما كانت على
حدود ما قبل الحرب.}%
}

%%K t12-15-1 p
15. --- ثم مرّت أمام أعيننا \VUgras{محطة البضائع} \textsuperscript{(71)}. وتحت
السقائف كان الموظفون يشحنون \VUgras{البضائع} \textsuperscript{(72)} المرسلة
بالقطار البطيء. و\VUgras{عمال المناورة} \textsuperscript{(76)} قرب \VUgras{خزان
الماء} \textsuperscript{(73)} كانوا ينقلون \VUgras{عربات الماشية} \textsuperscript{(75)}
و\VUgras{العربات المسطحة} \textsuperscript{(79)} ليؤلفوا \VUgras{قطار بضائع}
\textsuperscript{(80)}.

%%K t12-16-1 p
16. --- وعبرنا آخر \VUgras{مُحوّل} \textsuperscript{(78)} وكان عامل التحويل واقفًا
إلى جانبه؛ ومررنا بـ\VUgras{الإشارة} \textsuperscript{(86)} فاختفت المحطة في
البعد.

%%K t12-17-1 p
17. --- وسرعان ما خرجنا على \VUgras{الجسر الحديدي} \textsuperscript{(74)} الذي
يعبر \VUgras{النهر} \textsuperscript{(81)}. وبعد الجسر جرينا بمحاذاة النهر، وعليه
\VUgras{قاطرات} \textsuperscript{(82)} قوية تجرّ \VUgras{صنادل} \textsuperscript{(83)}
ثقيلة. ورأينا أيضًا \VUgras{سفينة ركاب} \textsuperscript{(84)} ترسو لتوّها عند
\VUgras{رصيف عائم} \textsuperscript{(85)}.

%%K t12-18-1 p
18. --- وبعد أن اجتزنا عدة محطات بسرعة، مررنا ببضعة \VUgras{أنفاق}
\textsuperscript{(87)} وتركنا خلفنا \VUgras{بيوتًا} \textsuperscript{(88)} لا تُحصى
لـ\VUgras{حرّاس المزلقانات}. وهدهدني تمايل القطار فغرقت في نوم عميق.

%%K t12-tit-4 sub
\VUsaut{3.0mm}
\VUpk{10.2pt}{40}{محطة دويتش-أفريكور.}
\VUsaut{3.0mm}

%%K t12-19-1 p
19. --- وأيقظني فجأةً صوت موظف يصيح: «دويتش-أفريكور!
\textsuperscript{(*)} الجميع يبدّلون!» كان ذلك تفتيش الجمرك وكل إجراءات الحدود
المملّة. واضطررت إلى ضبط ساعتي على \VUgras{ساعة} \textsuperscript{(89)} المحطة،
وكانت متقدمة خمسًا وخمسين دقيقة؛ وكنت بالكاد مستيقظًا فلم أكد أميّز
\VUgras{العقرب الكبير} \textsuperscript{(90)} من \VUgras{الصغير} \textsuperscript{(91)}؛
بل كان \VUgras{الميناء} \textsuperscript{(92)} نفسه مطموسًا بضباب الصباح.

%%K t12-20-1 p
20. --- وبينما كان موظفو الجمرك يجرون تفتيشهم، تثاءبت أتصفح
\VUgras{الملصقات} \textsuperscript{(93)} التي تزيّن جدران المحطة. وتناولت الفطور
لأقطع الوقت، وراجعت خريطة \VUgras{الشبكة} \textsuperscript{(94)} الألمانية لأرى كم
بقي بيني وبين ستراسبورغ، وعدت إلى عربتي وقد شدّ من أزري أملُ بلوغ غاية
سفري قريبًا.
\VUsaut{6.0mm}
\VUfilet{22mm}
